\documentclass[11pt]{article}

% parameters are fairly self-explanatory here, thankfully
\usepackage[width=7in,height=9.5in,top=0.75in,centering]{geometry}

% for math-ing
\usepackage{amsmath}
\usepackage{amssymb}


%%%%%% tables %%%%%%%%

\usepackage{array}
\renewcommand{\arraystretch}{1.5}
\setlength{\tabcolsep}{8pt}

%%%%%% plotting %%%%%%

\usepackage[dvipsnames]{xcolor}
\usepackage{tikz}
\usepackage{pgfplots}

\pgfplotsset{every axis/.append style={
                    axis x line=middle,    % put the x axis in the middle
                    axis y line=middle,    % put the y axis in the middle
                    axis line style={<->,color=black}, % arrows on the axis
                    xlabel={$x$},          % default put x on x-axis
                    ylabel={$y$},          % default put y on y-axis
            }}
            
 \pgfplotsset{compat=1.14} % sometimes, everything falls apart without this
 
 %\usepgfplotslibrary{fillbetween} %for shading between curves

%%%%%% commands %%%%%%

% exam heading; course name is first argument, exam information is second argument
\newcommand{\Heading}[2]{
    \fbox{
        {\Large\textbf{#1}}
        }
    \hfill
    \fbox{
        \begin{minipage}{0.2\textwidth}
            \begin{center}
            \textbf{#2}
            \end{center}
        \end{minipage}
    }
    \par\vspace{0.1in}}

% for being lazy while beautifying integrals
\newcommand{\dint}{\displaystyle\int}

% for being lazy while beautifying limits
\newcommand{\dlim}{\displaystyle\lim}

%%%%%% stuff %%%%%%

\usepackage{fancyhdr}

%\setcounter{page}{0} % first page is page 0

\parindent0pt % no indent

\fboxsep=0.25cm % little space in fbox border

%%%%%% BEGIN! %%%%%%%

\begin{document}

\pagestyle{fancy}

% record goals here to create sense of transparency

\fancyhead[c]{%
  \ifcase\value{page}
  % {\bf\color{ProcessBlue} F1: Lines \& Slope}% page=1
  % \or {\bf\color{RoyalBlue} L1: Limits and Continuity}% page = 2
  % \or {\bf\color{RoyalBlue} L2: Evaluating Limits Using Continuity and Forms}% page = 3
  % \or {\bf\color{ForestGreen} D1: Meaning of a Derivative}% page = 4
  \or {\bf\color{ProcessBlue} F2: Exponential and Logarithmic Functions}% page = 5
  \or {\bf\color{ForestGreen} D3: Derivatives of Basic Functions}
  \quad\textbf{+}\quad
  {\bf\color{ForestGreen} D2: Compute Derivatives}% page = 6
  \or {\bf\color{ForestGreen} D3: Derivatives of Basic Functions}
  \quad\textbf{+}\quad
  {\bf\color{ForestGreen} D2: Compute Derivatives}% page = 6
  \or {\bf\color{RedOrange} A1: Derivatives and Graphing}% page = 8
  \or
  \else
  % Empty chead here!
  \fi
}

% \thispagestyle{empty} % don't display that this is page 0

\Heading{MATH 1210}{Review Session 2\\Spring 2026\\
Section 930}

\vspace{0.1in}

\textbf{Name} \hrulefill
% \newpage%
%%%%%%%%%


\begin{enumerate}

%%%%%%%%%%%%%%%%%%%%%%%%%%%%%%%%%%%%%%%%%%%%%%%%%%%%%
%%%%%%%%%%% BEGIN PROBLEM F2 %%%%%%%%%%%%%%%%%%%%%%%%
%%%%%%%%%%%%%%%%%%%%%%%%%%%%%%%%%%%%%%%%%%%%%%%%%%%%%

\item[\textbf{F2.}]
For this question, you are not required to compute any powers of constants or logarithms.
For example, if your final answer includes $2^3$ or $\log_3(6)$, you can leave it as is.

\begin{enumerate}
\item
Simplify the expression below as much as possible so that the final result contains \textbf{only positive exponents}.
\begin{enumerate}
\item
$\left(\dfrac{\sqrt[3]{25} \ a^{-3/4} \ b^{-2/3}}{\sqrt{4} \ a^{-1/2} \ b^{2/3}}\right)^3$
\vfill
\item
$\left(\dfrac{49 \ m^{-1/2} \ n^{-3/5}}{7^2 \ m^{-1/4} \ n^{-1/5}}\right)^{11}$
\end{enumerate}
\vfill

\item Solve the following equation for $x$. 
\begin{enumerate}
\item
$5^{2x} + 2\cdot 5^x - 3 = 0$
\vfill
\item
$\log_{5}(x) - \log_{5}(x^3-7x^2) + \log_{5}(x - 7) = 3$
\end{enumerate}
~~\vfill
\end{enumerate}


%%%%%%%%%%%%%%%%%%%%%%%%%%%%%%%%%%%%%%%%%%%%%%%%%%%%%
%%%%%%%%%%% END PROBLEM F2 %%%%%%%%%%%%%%%%%%%%%%%%%%
%%%%%%%%%%%%%%%%%%%%%%%%%%%%%%%%%%%%%%%%%%%%%%%%%%%%%

\newpage

%%%%%%%%%%%%%%%%%%%%%%%%%%%%%%%%%%%%%%%%%%%%%%%%%%%%%
%%%%%%%%%%% BEGIN PROBLEM D2/D3 %%%%%%%%%%%%%%%%%%%%%%%%
%%%%%%%%%%%%%%%%%%%%%%%%%%%%%%%%%%%%%%%%%%%%%%%%%%%%%

\item[\textbf{D2/D3.}]
Write down the derivatives of the following functions.
No need to simplify.
\begin{enumerate}

\item $\displaystyle \frac{d}{dx}\left[ \frac{\ln(x)}{x^3} \right]$
\vfill
\item $\displaystyle \frac{d}{dx}\left[ \frac{x e^{x^2}}{x^2+1} \right]$
\vfill
\item $\displaystyle \frac{d}{dx}\left[ (x^3+1)e^{-2x}\right]$
\vfill
\end{enumerate}


%%%%%%%%%%%%%%%%%%%%%%%%%%%%%%%%%%%%%%%%%%%%%%%%%%%%%
%%%%%%%%%%% END PROBLEM D2/D3 %%%%%%%%%%%%%%%%%%%%%%%%%%
%%%%%%%%%%%%%%%%%%%%%%%%%%%%%%%%%%%%%%%%%%%%%%%%%%%%%

\newpage

%%%%%%%%%%%%%%%%%%%%%%%%%%%%%%%%%%%%%%%%%%%%%%%%%%%%%
%%%%%%%%%%% BEGIN PROBLEM D2/D3 %%%%%%%%%%%%%%%%%%%%%%%%
%%%%%%%%%%%%%%%%%%%%%%%%%%%%%%%%%%%%%%%%%%%%%%%%%%%%%

\item[\textbf{D2/D3.}]
Write down the derivative of the following function.
No need to simplify.
\[\displaystyle \frac{d}{dx}\left[ \ln\left(\frac{(x^2+3)e^{x}}{x} \right)\right]\]
\vfill


%%%%%%%%%%%%%%%%%%%%%%%%%%%%%%%%%%%%%%%%%%%%%%%%%%%%%
%%%%%%%%%%% END PROBLEM D2/D3 %%%%%%%%%%%%%%%%%%%%%%%%%%
%%%%%%%%%%%%%%%%%%%%%%%%%%%%%%%%%%%%%%%%%%%%%%%%%%%%%

\newpage

%%%%%%%%%%%%%%%%%%%%%%%%%%%%%%%%%%%%%%%%%%%%%%%%%%%%%
%%%%%%%%%%% BEGIN PROBLEM A1 %%%%%%%%%%%%%%%%%%%%%%%%
%%%%%%%%%%%%%%%%%%%%%%%%%%%%%%%%%%%%%%%%%%%%%%%%%%%%%

\item[\textbf{A1.}]
Consider the function $f(x)=\frac{x^2}{4} - 2x + \ln(x-1)+11$.
On what interval(s) is $f(x)$ increasing/decreasing?
On what interval(s) is $f(x)$ concave up/down?


%%%%%%%%%%%%%%%%%%%%%%%%%%%%%%%%%%%%%%%%%%%%%%%%%%%%%
%%%%%%%%%%% END PROBLEM A1 %%%%%%%%%%%%%%%%%%%%%%%%%%
%%%%%%%%%%%%%%%%%%%%%%%%%%%%%%%%%%%%%%%%%%%%%%%%%%%%%

\end{enumerate}

\end{document}